\documentclass[12pt, ukrainian]{article}
\usepackage[a4paper]{geometry}
\setlength\parindent{0mm}
\geometry{top=1.1cm,bottom=1.1cm,left=1.0cm,right=1.0cm}
\pagestyle{empty}
\usepackage[T2A]{fontenc}
\usepackage[utf8]{inputenc}
\usepackage[ukrainian]{babel}
\usepackage{graphicx}
\usepackage{caption}
\usepackage{caption}
\DeclareCaptionFormat{nocaption}{}
\captionsetup{format=nocaption,aboveskip=0pt,belowskip=0pt}
\usepackage[Export]{adjustbox} 
\adjustboxset{max size={0.9\linewidth}{0.9\paperheight}}
\usepackage{verbatim, listings, clrscode3e, fancyvrb}
\def\code#1{\Verb!#1!}
\begin{document}
%begin
\begin {large}

\textbf{01.12.2023 Контрольна робота, варіант  \No { 1 }}\par


{
%beginitem
1. Нехай int var=13; Один потік збільшив змінну var на 3, а інший збільшив її в три рази.

гонка даних (data race), стан гонки (race condition)

Виберіть всі правильні твердження.\par
\vspace{2mm}
( 1 ) Тут гарантовано є стан гонки.%good

( 2 ) Тут гарантовано є стан гонки, але гонки даних може не бути.%good

( 3 ) Тут гарантовано нема стану гонки, але може бути гонка даних.%bad

( 4 ) Тут гарантовано є стан гонки.%good
%enditem


\vspace{5mm}
%beginitem
2. Виберіть усі правильні твердження (якщо такі є).\par
\vspace{2mm}

( 1 ) Клас thread має засоби, щоб повернути результат виконання запущеного потоку.%bad

( 2 ) Розблоковувати незаблокований м’ютекс не можна.%good

( 3 ) Об“єкти класів promise та future можна використовувати тільки в багатопотокових програмах.%bad

( 4 ) Критична секція це неперервна ділянка коду програми.%bad
%enditem


\vspace{5mm}
%beginitem
3. Виберіть усі правильні твердження (якщо такі є).\par
\vspace{2mm}

( 1 ) Структура даних без блокувань зобов’язана використовувати атомарні змінні.%bad

( 2 ) Вільна від блокувань структура даних гарантує, що жоден потік не буде очікувати доступу до виконання операцій над цією структурою даних.%bad

( 3 ) Одночасно використовувати атомарних змінних більше, ніж регістрів процесора, у програмі не можна.%bad

( 4 ) Реалізація структури даних з блокуваннями може використовувати атомарні змінні.%good
%enditem


\vspace{5mm}
%beginitem
4. Що є точкою входу в новий потік, що його запущено засобами бібліотеки thread? Чи може батьківський потік отримати з нього значення-результат? Чи можна і якщо можна, то як, у новий потік передати посилання на дані батьківського потоку?

%enditem


\vspace{5mm}
%beginitem
5. Чи можна програмувати мовою С++ у стилі декларативного програмування? Відповідь пояснити.
%enditem
}
%end
\end {large}

\vspace{4mm}
%end
\newpage
%begin
\begin {large}

\textbf{01.12.2023 Контрольна робота, варіант  \No { 2 }}\par


{
%beginitem
1. Нехай int var=13; Один потік збільшив змінну var на 3, а інший збільшив її в три рази.

гонка даних (data race), стан гонки (race condition)

Виберіть всі правильні твердження.\par
\vspace{2mm}
( 1 ) Тут може не бути стану гонки.%bad

( 2 ) Тут гарантовано нема гонки даних, але може бути стан гонки.%bad

( 3 ) Тут гарантовано є стан гонки, але гонки даних може не бути.%good

( 4 ) Тут гарантовано нема ані гонки даних, ані стану гонки.%bad
%enditem


\vspace{5mm}
%beginitem
2. Виберіть усі правильні твердження (якщо такі є).\par
\vspace{2mm}

( 1 ) Секція є критичною, якщо під час її виконання не можна генерувати винятки.%bad

( 2 ) Потік в активному очікуванні споживає ресурс процесора.%good

( 3 ) Заблокований м’ютекс може розблокувати довільний потік.%bad

( 4 ) Тільки той потік, який обнулив значення семафора, може його збільшити.%bad
%enditem


\vspace{5mm}
%beginitem
3. Виберіть усі правильні твердження (якщо такі є).\par
\vspace{2mm}

( 1 ) Операція над вільною від блокувань структурою даних завжди виконається за обмежений час.%bad

( 2 ) За використання структурою даних рівно одного м’ютексу, взаємоблокування на її операціях неможливі.%good

( 3 ) Структура даних без блокувань гарантує завершення виконання довільної операції над цією структурою даних за обмежений час.%bad

( 4 ) Реалізації різних операцій над структурою даних повинні використовувати різні м’ютекси.%bad
%enditem


\vspace{5mm}
%beginitem
4. Чи завжди наслідком стану гонки (race condition) є невизначена поведінка програми? Відповідь пояснити.

%enditem


\vspace{5mm}
%beginitem
5. Що з наведеного переліку (якщо таке в переліку є):

\vspace{4mm}
\qquad нормальний порядок редукцій, ледачі обчислення, структурне програмування, анонімні функції

\vspace{4mm}
має найслабший логічний зв’язок з більшістю понять (відноситься до іншої парадигми та/або виривається з контексту чи не дуже сумісне)? Відповідь пояснити.

%enditem
}
%end
\end {large}

\vspace{4mm}
%end
\newpage
%begin
\begin {large}

\textbf{01.12.2023 Контрольна робота, варіант  \No { 3 }}\par


{
%beginitem
1. Нехай int var=13; Один потік збільшив змінну var на 3, а інший збільшив її в три рази.

гонка даних (data race), стан гонки (race condition)

Виберіть всі правильні твердження.\par
\vspace{2mm}
( 1 ) Тут гарантовано нема ані гонки даних, ані стану гонки.%bad

( 2 ) Тут може не бути гонки даних.%good

( 3 ) Тут гарантовано є гонка даних.%bad

( 4 ) Тут може не бути ані гонки даних, ані стану гонки.%bad
%enditem


\vspace{5mm}
%beginitem
2. Виберіть усі правильні твердження (якщо такі є).\par
\vspace{2mm}

( 1 ) Досвічений програміст за однопотоковою програмою (текстом мовою С++) завжди може вказати послідовність виконуваних обчислень на конкретних вхідних даних.%bad

( 2 ) Потік, запущений конструктором класу jthread, має обов’язково зупинитися, якщо пов’язаний об’єкт надішле команду зупинитися.%bad

( 3 ) Розблоковувати незаблокований м’ютекс можна.%bad

( 4 ) Розблокувати м’ютекс може тільки той потік, що його заблокував.%good
%enditem


\vspace{5mm}
%beginitem
3. Виберіть усі правильні твердження (якщо такі є).\par
\vspace{2mm}

( 1 ) Структура даних є структурою даних без блокувань, якщо її реалізація не використовує м’ютекси.%bad

( 2 ) Мова С++20 має стандартну бібліотеку конкурентних структур даних.%bad

( 3 ) Стандартні контейнери stl безпечно використовувати з кількох потоків.%bad

( 4 ) Структура даних з блокуваннями може бути такою, що використовує атомарні змінні.%good
%enditem


\vspace{5mm}
%beginitem
4. Що є наслідком захисту м’ютексами занадто великого фрагменту коду?

%enditem


\vspace{5mm}
%beginitem
5. Що з наступного не підтримує ані мова С++, ані її стандартна бібліотека: каррінг, виклик за необхідністю, замикання, ледачі обчислення, строгі (strict) функції? Відповідь пояснити.
%enditem
}
%end
\end {large}

\vspace{4mm}
%end
\newpage
%begin
\begin {large}

\textbf{01.12.2023 Контрольна робота, варіант  \No { 4 }}\par


{
%beginitem
1. Нехай int var=13; Один потік збільшив змінну var на 3, а інший збільшив її в три рази.

гонка даних (data race), стан гонки (race condition)

Виберіть всі правильні твердження.\par
\vspace{2mm}
( 1 ) Тут гарантовано є і стан гонки, і гонка даних.%bad

( 2 ) Тут гарантовано є і стан гонки, і гонка даних.%bad

( 3 ) Тут гарантовано є і стан гонки, і гонка даних.%bad

( 4 ) Тут гарантовано є стан гонки.%good
%enditem


\vspace{5mm}
%beginitem
2. Виберіть усі правильні твердження (якщо такі є).\par
\vspace{2mm}

( 1 ) Довільний потік може розблокувати заблокований м’ютекс.%bad

( 2 ) Щоб розблокувати заблокований \code{recursive\_mutex}, його завжди достатньо розблокувати один раз.%bad

( 3 ) Критичну секцію можна реалізувати за допомогою м’ютексу.%good

( 4 ) М’ютекс має бути записаний на початку критичної секції, яку він забезпечує.%bad
%enditem


\vspace{5mm}
%beginitem
3. Виберіть усі правильні твердження (якщо такі є).\par
\vspace{2mm}

( 1 ) Вільних від очікувань структур даних не існує.%bad

( 2 ) Структура даних може використовувати м’ютекси та атомарні змінні одночасно.%good

( 3 ) Мова С++20 має стандартну бібліотеку вільних від блокувань структур даних.%bad

( 4 ) Для уникнення гонки даних структура даних може використовувати не більше одного м’ютекса.%bad
%enditem


\vspace{5mm}
%beginitem
4. У наслідок чого можуть виникати мертві блокування (deadlock)? Як з ними можна боротися? Які є бібліотечні засоби, що частково вирішують проблему?

%enditem


\vspace{5mm}
%beginitem
5. Чи підходить мова Haskell для програмування в імперативному стилі? Відповідь пояснити.
%enditem
}
%end
\end {large}

\vspace{4mm}
%end
\newpage
%begin
\begin {large}

\textbf{01.12.2023 Контрольна робота, варіант  \No { 5 }}\par


{
%beginitem
1. Нехай int var=13; Один потік збільшив змінну var на 3, а інший збільшив її в три рази.

гонка даних (data race), стан гонки (race condition)

Виберіть всі правильні твердження.\par
\vspace{2mm}
( 1 ) Тут може не бути стану гонки.%bad

( 2 ) Тут може не бути гонки даних.%good

( 3 ) Тут може не бути стану гонки.%bad

( 4 ) Тут гарантовано нема гонки даних, але може бути стан гонки.%bad
%enditem


\vspace{5mm}
%beginitem
2. Виберіть усі правильні твердження (якщо такі є).\par
\vspace{2mm}

( 1 ) Текст однопотокової програми мовою С++ однозначно визначає послідовність виконуваних обчислень.%bad

( 2 ) Значення, що є результатом виконання функції, запущеної через async, зберігається безпосередньо в об’єкті ф’ючерса (future).%bad

( 3 ) Критична секція може складатися з кількох ділянок коду, які можуть розташовуватись де завгодно (не обов’язково зв’язний фрагмент коду).%good

( 4 ) Об’єкт future можна скопіювати та роздати копії різним потокам, щоб вони змогли побачити результат.%bad
%enditem


\vspace{5mm}
%beginitem
3. Виберіть усі правильні твердження (якщо такі є).\par
\vspace{2mm}

( 1 ) Взаємоблокування потоків може виникнути тільки за використання м’ютексів.%bad

( 2 ) Операція над вільною від блокувань структурою даних може виконуватись скільки завгодно довго.%good

( 3 ) Структура даних без блокувань та вільна від блокувань структура даних --- це одне те саме.%bad

( 4 ) Операція над структурою даних без блокувань завжди виконається за обмежений час.%bad
%enditem


\vspace{5mm}
%beginitem
4. Які проблеми можуть виникнути за concurrent-доступу до даних, що зберігаються в структурі даних з операціями завантажити дані та ітератором, що забезпечує її обхід? Як їх можна уникнути?

%enditem


\vspace{5mm}
%beginitem
5. Що з наведеного переліку (якщо таке в переліку є):

\vspace{4mm}
\qquad послідовність виконання інструкцій, прозорість за посиланнями, ледачі обчислення, анонімні функції

\vspace{4mm}
має найслабший логічний зв’язок з більшістю понять (відноситься до іншої парадигми та/або виривається з контексту чи не дуже сумісне)? Відповідь пояснити.

%enditem
}
%end
\end {large}

\vspace{4mm}
%end
\newpage
%begin
\begin {large}

\textbf{01.12.2023 Контрольна робота, варіант  \No { 6 }}\par


{
%beginitem
1. Нехай int var=13; Один потік збільшив змінну var на 3, а інший збільшив її в три рази.

гонка даних (data race), стан гонки (race condition)

Виберіть всі правильні твердження.\par
\vspace{2mm}
( 1 ) Тут може не бути ані гонки даних, ані стану гонки.%bad

( 2 ) Тут гарантовано є гонка даних.%bad

( 3 ) Тут гарантовано нема ані гонки даних, ані стану гонки.%bad

( 4 ) Тут гарантовано нема стану гонки, але може бути гонка даних.%bad
%enditem


\vspace{5mm}
%beginitem
2. Виберіть усі правильні твердження (якщо такі є).\par
\vspace{2mm}

( 1 ) Тільки від компілятора залежить, чи запустить функція async виконання в новому потоці.%bad

( 2 ) У програмі може бути не більше однієї критичної секції.%bad

( 3 ) Якщо два обчислення програми на С++ невдало, але синхронізовані, то на них може виникнути гонка даних.%good

( 4 ) Функція async запускає виконання в новому потоці.%bad
%enditem


\vspace{5mm}
%beginitem
3. Виберіть усі правильні твердження (якщо такі є).\par
\vspace{2mm}

( 1 ) Структура даних без блокувань гарантує, що жоден потік не буде очікувати доступу до виконання операцій над цією структурою даних.%bad

( 2 ) За використання атомарних змінних структура даних є структурою даних без блокувань.%bad

( 3 ) Стандартні контейнери stl безпечно використовувати з кількох потоків.%bad

( 4 ) Кожна вільна від очікувань структура даних є структурою даних, вільною від блокувань.%good
%enditem


\vspace{5mm}
%beginitem
4. Що з точки зору С++20 можна сказати про послідовність обчислень підвиразів у виразі f(R(13), new P(b+c), Q(a))) ? Вважати всі імена коректно означеними.

%enditem


\vspace{5mm}
%beginitem
5. Чи можна програмувати мовою Java у стилі декларативного програмування? Відповідь пояснити.
%enditem
}
%end
\end {large}

\vspace{4mm}
%end
\newpage
%begin
\begin {large}

\textbf{01.12.2023 Контрольна робота, варіант  \No { 7 }}\par


{
%beginitem
1. Нехай int var=13; Один потік збільшив змінну var на 3, а інший збільшив її в три рази.

гонка даних (data race), стан гонки (race condition)

Виберіть всі правильні твердження.\par
\vspace{2mm}
( 1 ) Тут гарантовано є гонка даних.%bad

( 2 ) Тут гарантовано нема гонки даних, але може бути стан гонки.%bad

( 3 ) Тут гарантовано є стан гонки, але гонки даних може не бути.%good

( 4 ) Тут гарантовано є стан гонки.%good
%enditem


\vspace{5mm}
%beginitem
2. Виберіть усі правильні твердження (якщо такі є).\par
\vspace{2mm}

( 1 ) У стані гонки програма відчайдушно намагається захопити ресурси процесора.%bad

( 2 ) Коли програма спить, то вона не виконується.%bad

( 3 ) У синтаксично коректній програмі мовою С++ стан гонки (race condition) може виникати.%good

( 4 ) Значення, що є результатом виконання функції, запущеної через async, зберігається безпосередньо в об’єкті проміза (promise).%bad
%enditem


\vspace{5mm}
%beginitem
3. Виберіть усі правильні твердження (якщо такі є).\par
\vspace{2mm}

( 1 ) Структура даних без блокувань гарантує завершення виконання довільної операції над цією структурою даних за обмежений час.%bad

( 2 ) Вільна від блокувань структура даних гарантує, що жоден потік не буде очікувати доступу до виконання операцій над цією структурою даних.%bad

( 3 ) Реалізація структури даних без блокувань не може використовувати м’ютекси.%good

( 4 ) Операція над вільною від блокувань структурою даних завжди виконається за обмежений час.%bad
%enditem


\vspace{5mm}
%beginitem
4. Які є засоби С++20 та прийоми, щоб забезпечити безпечне створення глобальних даних?

%enditem


\vspace{5mm}
%beginitem
5. Що з наведеного переліку (якщо таке в переліку є):

\vspace{4mm}
\qquad апплікативний порядок редукцій, строгі (strict) функції, каррінг, виклик за значенням

\vspace{4mm}
має найслабший логічний зв’язок з більшістю понять (відноситься до іншої парадигми та/або виривається з контексту чи не дуже сумісне)? Відповідь пояснити.

%enditem
}
%end
\end {large}

\vspace{4mm}
%end
\newpage
%begin
\begin {large}

\textbf{01.12.2023 Контрольна робота, варіант  \No { 8 }}\par


{
%beginitem
1. Нехай int var=13; Один потік збільшив змінну var на 3, а інший збільшив її в три рази.

гонка даних (data race), стан гонки (race condition)

Виберіть всі правильні твердження.\par
\vspace{2mm}
( 1 ) Тут гарантовано нема стану гонки, але може бути гонка даних.%bad

( 2 ) Тут може не бути гонки даних.%good

( 3 ) Тут може не бути стану гонки.%bad

( 4 ) Тут гарантовано є стан гонки, але гонки даних може не бути.%good
%enditem


\vspace{5mm}
%beginitem
2. Виберіть усі правильні твердження (якщо такі є).\par
\vspace{2mm}

( 1 ) Функція \code{sleep\_for} може змусити інший потік заснути.%bad

( 2 ) Блокувати заблокований м’ютекс  \code{mutex} іноді можна.%good

( 3 ) Коли програма спить, то сплять усі її потоки.%bad

( 4 ) Не можна блокувати заблокований м’ютекс mutex.%bad
%enditem


\vspace{5mm}
%beginitem
3. Виберіть усі правильні твердження (якщо такі є).\par
\vspace{2mm}

( 1 ) Реалізації різних операцій над структурою даних повинні використовувати різні м’ютекси.%bad

( 2 ) За використання атомарних змінних структура даних є структурою даних без блокувань.%bad

( 3 ) Мова С++20 має стандартну бібліотеку вільних від блокувань структур даних.%bad

( 4 ) Якщо весь код структури даних утворює одну критичну секцію, то така структура даних може використовуватись в конкурентній програмі, але по-справжньому конкурентною не є.%good
%enditem


\vspace{5mm}
%beginitem
4. Що з точки зору С++20 можна сказати про послідовність обчислень підвиразів у виразі f(a+b)+g(c+d) ? Вважати всі імена коректно означеними.

%enditem


\vspace{5mm}
%beginitem
5. Чи можна програмувати мовою Java у стилі структурного програмування? Відповідь пояснити.
%enditem
}
%end
\end {large}

\vspace{4mm}
%end
\newpage
%begin
\begin {large}

\textbf{01.12.2023 Контрольна робота, варіант  \No { 9 }}\par


{
%beginitem
1. Нехай int var=13; Один потік збільшив змінну var на 3, а інший збільшив її в три рази.

гонка даних (data race), стан гонки (race condition)

Виберіть всі правильні твердження.\par
\vspace{2mm}
( 1 ) Тут може не бути ані гонки даних, ані стану гонки.%bad

( 2 ) Тут гарантовано нема стану гонки, але може бути гонка даних.%bad

( 3 ) Тут гарантовано є і стан гонки, і гонка даних.%bad

( 4 ) Тут гарантовано нема гонки даних, але може бути стан гонки.%bad
%enditem


\vspace{5mm}
%beginitem
2. Виберіть усі правильні твердження (якщо такі є).\par
\vspace{2mm}

( 1 ) Об’єкт promise можна скопіювати та роздати копії різним потокам, щоб кожен з них міг передати результат своєї роботи.%bad

( 2 ) Виклик \code{this\_thread::sleep\_for(5s)} змушує потік проспати рівно 5 секунд.%bad

( 3 ) Коли програма спить, то сплять усі її потоки.%bad

( 4 ) Коли потік спить, він не споживає ресурс процесора.%good
%enditem


\vspace{5mm}
%beginitem
3. Виберіть усі правильні твердження (якщо такі є).\par
\vspace{2mm}

( 1 ) Вільних від очікувань структур даних не існує.%bad

( 2 ) Мова С++20 має стандартну бібліотеку конкурентних структур даних.%bad

( 3 ) Взаємоблокування потоків може виникнути тільки за використання м’ютексів.%bad

( 4 ) Реалізація структури даних без блокувань не може використовувати м’ютекси.%good
%enditem


\vspace{5mm}
%beginitem
4. Які бібліотечні засоби С++20 та як здатні забезпечувати захист глобальних даних від стану гонки, якщо може виконуватись довільна кількість операцій доступу (читання та запису) до цих даних?

%enditem


\vspace{5mm}
%beginitem
5. Що з наведеного переліку (якщо таке в переліку є):

\vspace{4mm}
\qquad енергійні обчислення, виклик за необхідністю, прозорість за посиланнями, каррінг

\vspace{4mm}
має найслабший логічний зв’язок з більшістю понять (відноситься до іншої парадигми та/або виривається з контексту чи не дуже сумісне)? Відповідь пояснити.

%enditem
}
%end
\end {large}

\vspace{4mm}
%end
\newpage
%begin
\begin {large}

\textbf{01.12.2023 Контрольна робота, варіант  \No { 10 }}\par


{
%beginitem
1. Нехай int var=13; Один потік збільшив змінну var на 3, а інший збільшив її в три рази.

гонка даних (data race), стан гонки (race condition)

Виберіть всі правильні твердження.\par
\vspace{2mm}
( 1 ) Тут може не бути гонки даних.%good

( 2 ) Тут може не бути ані гонки даних, ані стану гонки.%bad

( 3 ) Тут гарантовано нема ані гонки даних, ані стану гонки.%bad

( 4 ) Тут гарантовано є гонка даних.%bad
%enditem


\vspace{5mm}
%beginitem
2. Виберіть усі правильні твердження (якщо такі є).\par
\vspace{2mm}

( 1 ) Коли програма спить, то вона не виконується.%bad

( 2 ) Об“єкти класів promise та future можна використовувати тільки в багатопотокових програмах.%bad

( 3 ) Виклик функції \code{this\_thread::sleep\_for(5s)} змушує потік, з якого його здійснили заснути якнайменш на 5 секунд.%good

( 4 ) Потік, запущений конструктором класу jthread, має обов’язково зупинитися, якщо пов’язаний об’єкт надішле команду зупинитися.%bad
%enditem


\vspace{5mm}
%beginitem
3. Виберіть усі правильні твердження (якщо такі є).\par
\vspace{2mm}

( 1 ) Реалізація структури даних з блокуваннями може використовувати атомарні змінні.%good

( 2 ) Структура даних без блокувань та вільна від блокувань структура даних --- це одне те саме.%bad

( 3 ) Операція над структурою даних без блокувань завжди виконається за обмежений час.%bad

( 4 ) Структура даних без блокувань гарантує, що жоден потік не буде очікувати доступу до виконання операцій над цією структурою даних.%bad
%enditem


\vspace{5mm}
%beginitem
4. Які бібліотечні засоби С++20 дозволяють не тільки запустити код в іншому потоці, а ще без особливих танців з бубнами дізнатись результат його виконання?

%enditem


\vspace{5mm}
%beginitem
5. Чи можна одну ту саму мову програмування використовувати для програмування у функціональному, імперативному та об’єктно-орієнтованому стилі одночасно? Відповідь пояснити.
%enditem
}
%end
\end {large}

\vspace{4mm}
%end
\newpage
%begin
\begin {large}

\textbf{01.12.2023 Контрольна робота, варіант  \No { 11 }}\par


{
%beginitem
1. Нехай int var=13; Один потік збільшив змінну var на 3, а інший збільшив її в три рази.

гонка даних (data race), стан гонки (race condition)

Виберіть всі правильні твердження.\par
\vspace{2mm}
( 1 ) Тут гарантовано нема стану гонки, але може бути гонка даних.%bad

( 2 ) Тут гарантовано є гонка даних.%bad

( 3 ) Тут гарантовано є стан гонки, але гонки даних може не бути.%good

( 4 ) Тут гарантовано нема ані гонки даних, ані стану гонки.%bad
%enditem


\vspace{5mm}
%beginitem
2. Виберіть усі правильні твердження (якщо такі є).\par
\vspace{2mm}

( 1 ) Чи запустить функція  async виконання в новому потоці можна задавати політикою.%good

( 2 ) Тільки той потік, який обнулив значення семафора, може його збільшити.%bad

( 3 ) Щоб розблокувати заблокований \code{recursive\_mutex}, його завжди достатньо розблокувати один раз.%bad

( 4 ) Тільки від компілятора залежить, чи запустить функція async виконання в новому потоці.%bad
%enditem


\vspace{5mm}
%beginitem
3. Виберіть усі правильні твердження (якщо такі є).\par
\vspace{2mm}

( 1 ) Операція над вільною від блокувань структурою даних може виконуватись скільки завгодно довго.%good

( 2 ) Структура даних є структурою даних без блокувань, якщо її реалізація не використовує м’ютекси.%bad

( 3 ) Одночасно використовувати атомарних змінних більше, ніж регістрів процесора, у програмі не можна.%bad

( 4 ) Структура даних без блокувань зобов’язана використовувати атомарні змінні.%bad
%enditem


\vspace{5mm}
%beginitem
4. Чи можна обмежити доступ потоку до глобальних змінних? А якщо дуже треба? Відповідь пояснити.

%enditem


\vspace{5mm}
%beginitem
5. Чи можна програмувати мовою С++ у стилі структурного програмування? Відповідь пояснити.
%enditem
}
%end
\end {large}

\vspace{4mm}
%end
\newpage
%begin
\begin {large}

\textbf{01.12.2023 Контрольна робота, варіант  \No { 12 }}\par


{
%beginitem
1. Нехай int var=13; Один потік збільшив змінну var на 3, а інший збільшив її в три рази.

гонка даних (data race), стан гонки (race condition)

Виберіть всі правильні твердження.\par
\vspace{2mm}
( 1 ) Тут може не бути ані гонки даних, ані стану гонки.%bad

( 2 ) Тут гарантовано є і стан гонки, і гонка даних.%bad

( 3 ) Тут гарантовано нема гонки даних, але може бути стан гонки.%bad

( 4 ) Тут може не бути гонки даних.%good
%enditem


\vspace{5mm}
%beginitem
2. Виберіть усі правильні твердження (якщо такі є).\par
\vspace{2mm}

( 1 ) Досвічений програміст за однопотоковою програмою (текстом мовою С++) завжди може вказати послідовність виконуваних обчислень на конкретних вхідних даних.%bad

( 2 ) Щоб розблокувати заблокований \code{shared\_mutex}, його завжди достатньо розблокувати один раз.%good

( 3 ) У програмі може бути не більше однієї критичної секції.%bad

( 4 ) Функція \code{sleep\_for} може змусити інший потік заснути.%bad
%enditem


\vspace{5mm}
%beginitem
3. Виберіть усі правильні твердження (якщо такі є).\par
\vspace{2mm}

( 1 ) Кожна вільна від очікувань структура даних є структурою даних, вільною від блокувань.%good

( 2 ) Для уникнення гонки даних структура даних може використовувати не більше одного м’ютекса.%bad

( 3 ) Вільних від очікувань структур даних не існує.%bad

( 4 ) Вільна від блокувань структура даних гарантує, що жоден потік не буде очікувати доступу до виконання операцій над цією структурою даних.%bad
%enditem


\vspace{5mm}
%beginitem
4. Чи може потік мати доступ до автоматичної пам’яті іншого потоку? Відповідь пояснити.

%enditem


\vspace{5mm}
%beginitem
5. Що з наведеного переліку (якщо таке в переліку є):

\vspace{4mm}
\qquad функціональне програмування, чисті функції, каррінг, рекурсія, прозорість за посиланнями

\vspace{4mm}
має найслабший логічний зв’язок з більшістю понять (відноситься до іншої парадигми та/або виривається з контексту чи не дуже сумісне)? Відповідь пояснити.

%enditem
}
%end
\end {large}

\vspace{4mm}
%end
\newpage
%begin
\begin {large}

\textbf{01.12.2023 Контрольна робота, варіант  \No { 13 }}\par


{
%beginitem
1. Нехай int var=13; Один потік збільшив змінну var на 3, а інший збільшив її в три рази.

гонка даних (data race), стан гонки (race condition)

Виберіть всі правильні твердження.\par
\vspace{2mm}
( 1 ) Тут гарантовано нема стану гонки, але може бути гонка даних.%bad

( 2 ) Тут гарантовано є і стан гонки, і гонка даних.%bad

( 3 ) Тут гарантовано є стан гонки, але гонки даних може не бути.%good

( 4 ) Тут може не бути стану гонки.%bad
%enditem


\vspace{5mm}
%beginitem
2. Виберіть усі правильні твердження (якщо такі є).\par
\vspace{2mm}

( 1 ) Не можна блокувати заблокований м’ютекс mutex.%bad

( 2 ) М’ютекс має бути записаний на початку критичної секції, яку він забезпечує.%bad

( 3 ) Виклик \code{this\_thread::sleep\_for(5s)} змушує потік проспати рівно 5 секунд.%bad

( 4 ) Клас thread не має засобів, щоб повернути результат виконання запущеного потоку.%good
%enditem


\vspace{5mm}
%beginitem
3. Виберіть усі правильні твердження (якщо такі є).\par
\vspace{2mm}

( 1 ) За використання структурою даних рівно одного м’ютексу, взаємоблокування на її операціях неможливі.%good

( 2 ) Мова С++20 має стандартну бібліотеку конкурентних структур даних.%bad

( 3 ) Реалізації різних операцій над структурою даних повинні використовувати різні м’ютекси.%bad

( 4 ) Операція над вільною від блокувань структурою даних завжди виконається за обмежений час.%bad
%enditem


\vspace{5mm}
%beginitem
4. Чи мають відношення обіцянки (promise) до сопрограм? Відповідь пояснити.

%enditem


\vspace{5mm}
%beginitem
5. Чи підходить мова Haskell для програмування в імперативному стилі? Відповідь пояснити.
%enditem
}
%end
\end {large}

\vspace{4mm}
%end
\newpage
%begin
\begin {large}

\textbf{01.12.2023 Контрольна робота, варіант  \No { 14 }}\par


{
%beginitem
1. Нехай int var=13; Один потік збільшив змінну var на 3, а інший збільшив її в три рази.

гонка даних (data race), стан гонки (race condition)

Виберіть всі правильні твердження.\par
\vspace{2mm}
( 1 ) Тут гарантовано є стан гонки.%good

( 2 ) Тут гарантовано є стан гонки.%good

( 3 ) Тут може не бути ані гонки даних, ані стану гонки.%bad

( 4 ) Тут гарантовано є гонка даних.%bad
%enditem


\vspace{5mm}
%beginitem
2. Виберіть усі правильні твердження (якщо такі є).\par
\vspace{2mm}

( 1 ) Довільний потік може розблокувати заблокований м’ютекс.%bad

( 2 ) За різних запусків програми на одних тих самих вхідних даних може виникати хронологічно різна послідовність обчислень.%good

( 3 ) Клас thread має засоби, щоб повернути результат виконання запущеного потоку.%bad

( 4 ) Об’єкт future можна скопіювати та роздати копії різним потокам, щоб вони змогли побачити результат.%bad
%enditem


\vspace{5mm}
%beginitem
3. Виберіть усі правильні твердження (якщо такі є).\par
\vspace{2mm}

( 1 ) Структура даних без блокувань гарантує, що жоден потік не буде очікувати доступу до виконання операцій над цією структурою даних.%bad

( 2 ) Мова С++20 має стандартну бібліотеку вільних від блокувань структур даних.%bad

( 3 ) Для уникнення гонки даних структура даних може використовувати не більше одного м’ютекса.%bad

( 4 ) Структура даних може використовувати м’ютекси та атомарні змінні одночасно.%good
%enditem


\vspace{5mm}
%beginitem
4. Які бібліотечні засоби С++20 можна використати, що забезпечити передачу результатів обчислення від одного потоку до іншого? Яка схема їх використання?

%enditem


\vspace{5mm}
%beginitem
5. Чи можна програмувати мовою С++ у стилі структурного програмування? Відповідь пояснити.
%enditem
}
%end
\end {large}

\vspace{4mm}
%end
\newpage
%begin
\begin {large}

\textbf{01.12.2023 Контрольна робота, варіант  \No { 15 }}\par


{
%beginitem
1. Нехай int var=13; Один потік збільшив змінну var на 3, а інший збільшив її в три рази.

гонка даних (data race), стан гонки (race condition)

Виберіть всі правильні твердження.\par
\vspace{2mm}
( 1 ) Тут гарантовано нема гонки даних, але може бути стан гонки.%bad

( 2 ) Тут гарантовано нема гонки даних, але може бути стан гонки.%bad

( 3 ) Тут може не бути стану гонки.%bad

( 4 ) Тут гарантовано є гонка даних.%bad
%enditem


\vspace{5mm}
%beginitem
2. Виберіть усі правильні твердження (якщо такі є).\par
\vspace{2mm}

( 1 ) Порядок обчислень у потоці для одних тих самих вхідних даних може залежати від використовуваного компілятора.%good

( 2 ) Текст однопотокової програми мовою С++ однозначно визначає послідовність виконуваних обчислень.%bad

( 3 ) Значення, що є результатом виконання функції, запущеної через async, зберігається безпосередньо в об’єкті ф’ючерса (future).%bad

( 4 ) Секція є критичною, якщо під час її виконання не можна генерувати винятки.%bad
%enditem


\vspace{5mm}
%beginitem
3. Виберіть усі правильні твердження (якщо такі є).\par
\vspace{2mm}

( 1 ) Стандартні контейнери stl безпечно використовувати з кількох потоків.%bad

( 2 ) Операція над структурою даних без блокувань завжди виконається за обмежений час.%bad

( 3 ) Взаємоблокування потоків може виникнути тільки за використання м’ютексів.%bad

( 4 ) Структура даних з блокуваннями може бути такою, що використовує атомарні змінні.%good
%enditem


\vspace{5mm}
%beginitem
4. Як найкраще (з максимальним збереженням конкурентності) забезпечити захист глобальних даних від стану гонки, якщо над ними переважно виконуються операції читання і подекуди операції запису?

%enditem


\vspace{5mm}
%beginitem
5. Чи можна одну ту саму мову програмування використовувати для програмування у функціональному, імперативному та об’єктно-орієнтованому стилі одночасно? Відповідь пояснити.
%enditem
}
%end
\end {large}

\vspace{4mm}
%end
\newpage
%begin
\begin {large}

\textbf{01.12.2023 Контрольна робота, варіант  \No { 16 }}\par


{
%beginitem
1. Нехай int var=13; Один потік збільшив змінну var на 3, а інший збільшив її в три рази.

гонка даних (data race), стан гонки (race condition)

Виберіть всі правильні твердження.\par
\vspace{2mm}
( 1 ) Тут гарантовано нема ані гонки даних, ані стану гонки.%bad

( 2 ) Тут гарантовано є стан гонки, але гонки даних може не бути.%good

( 3 ) Тут може не бути гонки даних.%good

( 4 ) Тут гарантовано є стан гонки.%good
%enditem


\vspace{5mm}
%beginitem
2. Виберіть усі правильні твердження (якщо такі є).\par
\vspace{2mm}

( 1 ) Функція async запускає виконання в новому потоці.%bad

( 2 ) У стані гонки програма відчайдушно намагається захопити ресурси процесора.%bad

( 3 ) Об’єкт promise можна скопіювати та роздати копії різним потокам, щоб кожен з них міг передати результат своєї роботи.%bad

( 4 ) За різних запусків програми на одних тих самих вхідних даних може виникати хронологічно різна послідовність обчислень.%good
%enditem


\vspace{5mm}
%beginitem
3. Виберіть усі правильні твердження (якщо такі є).\par
\vspace{2mm}

( 1 ) Структура даних без блокувань та вільна від блокувань структура даних --- це одне те саме.%bad

( 2 ) Якщо весь код структури даних утворює одну критичну секцію, то така структура даних може використовуватись в конкурентній програмі, але по-справжньому конкурентною не є.%good

( 3 ) Структура даних без блокувань гарантує завершення виконання довільної операції над цією структурою даних за обмежений час.%bad

( 4 ) Одночасно використовувати атомарних змінних більше, ніж регістрів процесора, у програмі не можна.%bad
%enditem


\vspace{5mm}
%beginitem
4. Які бібліотечні засоби С++20 та за якою схемою (кратко) здатні забезпечувати синхронізацію потоків?

%enditem


\vspace{5mm}
%beginitem
5. Чи можна програмувати мовою Java у стилі структурного програмування? Відповідь пояснити.
%enditem
}
%end
\end {large}

\vspace{4mm}
%end
\newpage
%begin
\begin {large}

\textbf{01.12.2023 Контрольна робота, варіант  \No { 17 }}\par


{
%beginitem
1. Нехай int var=13; Один потік збільшив змінну var на 3, а інший збільшив її в три рази.

гонка даних (data race), стан гонки (race condition)

Виберіть всі правильні твердження.\par
\vspace{2mm}
( 1 ) Тут гарантовано нема стану гонки, але може бути гонка даних.%bad

( 2 ) Тут гарантовано є і стан гонки, і гонка даних.%bad

( 3 ) Тут гарантовано є і стан гонки, і гонка даних.%bad

( 4 ) Тут гарантовано нема ані гонки даних, ані стану гонки.%bad
%enditem


\vspace{5mm}
%beginitem
2. Виберіть усі правильні твердження (якщо такі є).\par
\vspace{2mm}

( 1 ) Заблокований м’ютекс може розблокувати довільний потік.%bad

( 2 ) Розблокувати м’ютекс може тільки той потік, що його заблокував.%good

( 3 ) Розблоковувати незаблокований м’ютекс можна.%bad

( 4 ) Значення, що є результатом виконання функції, запущеної через async, зберігається безпосередньо в об’єкті проміза (promise).%bad
%enditem


\vspace{5mm}
%beginitem
3. Виберіть усі правильні твердження (якщо такі є).\par
\vspace{2mm}

( 1 ) За використання атомарних змінних структура даних є структурою даних без блокувань.%bad

( 2 ) Реалізація структури даних без блокувань не може використовувати м’ютекси.%good

( 3 ) Структура даних без блокувань зобов’язана використовувати атомарні змінні.%bad

( 4 ) Структура даних є структурою даних без блокувань, якщо її реалізація не використовує м’ютекси.%bad
%enditem


\vspace{5mm}
%beginitem
4. Фьючерси та обіцянки (promise): як за допомогою них можна організувати асинхронне виконання частин коду?

%enditem


\vspace{5mm}
%beginitem
5. Що з наступного не підтримує ані мова С++, ані її стандартна бібліотека: каррінг, виклик за необхідністю, замикання, ледачі обчислення, строгі (strict) функції? Відповідь пояснити.
%enditem
}
%end
\end {large}

\vspace{4mm}
%end
\newpage
%begin
\begin {large}

\textbf{01.12.2023 Контрольна робота, варіант  \No { 18 }}\par


{
%beginitem
1. Нехай int var=13; Один потік збільшив змінну var на 3, а інший збільшив її в три рази.

гонка даних (data race), стан гонки (race condition)

Виберіть всі правильні твердження.\par
\vspace{2mm}
( 1 ) Тут може не бути стану гонки.%bad

( 2 ) Тут може не бути гонки даних.%good

( 3 ) Тут гарантовано нема гонки даних, але може бути стан гонки.%bad

( 4 ) Тут гарантовано є стан гонки, але гонки даних може не бути.%good
%enditem


\vspace{5mm}
%beginitem
2. Виберіть усі правильні твердження (якщо такі є).\par
\vspace{2mm}

( 1 ) Коли потік спить, він не споживає ресурс процесора.%good

( 2 ) Критична секція це неперервна ділянка коду програми.%bad

( 3 ) Об“єкти класів promise та future можна використовувати тільки в багатопотокових програмах.%bad

( 4 ) Не можна блокувати заблокований м’ютекс mutex.%bad
%enditem


\vspace{5mm}
%beginitem
3. Виберіть усі правильні твердження (якщо такі є).\par
\vspace{2mm}

( 1 ) Структура даних є структурою даних без блокувань, якщо її реалізація не використовує м’ютекси.%bad

( 2 ) За використання атомарних змінних структура даних є структурою даних без блокувань.%bad

( 3 ) Структура даних без блокувань зобов’язана використовувати атомарні змінні.%bad

( 4 ) Структура даних може використовувати м’ютекси та атомарні змінні одночасно.%good
%enditem


\vspace{5mm}
%beginitem
4. Як співвідносяться потоки операційної системи з потоками-об’єктами thread? Який з методів join, detach є блокуючим?

%enditem


\vspace{5mm}
%beginitem
5. Чи можна програмувати мовою С++ у стилі декларативного програмування? Відповідь пояснити.
%enditem
}
%end
\end {large}

\vspace{4mm}
%end
\newpage
%begin
\begin {large}

\textbf{01.12.2023 Контрольна робота, варіант  \No { 19 }}\par


{
%beginitem
1. Нехай int var=13; Один потік збільшив змінну var на 3, а інший збільшив її в три рази.

гонка даних (data race), стан гонки (race condition)

Виберіть всі правильні твердження.\par
\vspace{2mm}
( 1 ) Тут гарантовано є стан гонки.%good

( 2 ) Тут гарантовано нема стану гонки, але може бути гонка даних.%bad

( 3 ) Тут може не бути ані гонки даних, ані стану гонки.%bad

( 4 ) Тут може не бути стану гонки.%bad
%enditem


\vspace{5mm}
%beginitem
2. Виберіть усі правильні твердження (якщо такі є).\par
\vspace{2mm}

( 1 ) Критична секція може складатися з кількох ділянок коду, які можуть розташовуватись де завгодно (не обов’язково зв’язний фрагмент коду).%good

( 2 ) Заблокований м’ютекс може розблокувати довільний потік.%bad

( 3 ) У стані гонки програма відчайдушно намагається захопити ресурси процесора.%bad

( 4 ) Коли програма спить, то сплять усі її потоки.%bad
%enditem


\vspace{5mm}
%beginitem
3. Виберіть усі правильні твердження (якщо такі є).\par
\vspace{2mm}

( 1 ) Якщо весь код структури даних утворює одну критичну секцію, то така структура даних може використовуватись в конкурентній програмі, але по-справжньому конкурентною не є.%good

( 2 ) Для уникнення гонки даних структура даних може використовувати не більше одного м’ютекса.%bad

( 3 ) Стандартні контейнери stl безпечно використовувати з кількох потоків.%bad

( 4 ) Операція над вільною від блокувань структурою даних завжди виконається за обмежений час.%bad
%enditem


\vspace{5mm}
%beginitem
4. Що таке недетерміновано впорядковані (indeterminately sequenced) обчислення? Чи є їх наявність помилкою у програмі?  Відповідь пояснити.

%enditem


\vspace{5mm}
%beginitem
5. Що з наведеного переліку (якщо таке в переліку є):

\vspace{4mm}
\qquad взаємодія об’єктів, чисті функції, абстракція даних, побічний ефект

\vspace{4mm}
має найслабший логічний зв’язок з більшістю понять (відноситься до іншої парадигми та/або виривається з контексту чи не дуже сумісне)? Відповідь пояснити.

%enditem
}
%end
\end {large}

\vspace{4mm}
%end
\newpage
%begin
\begin {large}

\textbf{01.12.2023 Контрольна робота, варіант  \No { 20 }}\par


{
%beginitem
1. Нехай int var=13; Один потік збільшив змінну var на 3, а інший збільшив її в три рази.

гонка даних (data race), стан гонки (race condition)

Виберіть всі правильні твердження.\par
\vspace{2mm}
( 1 ) Тут гарантовано є гонка даних.%bad

( 2 ) Тут може не бути ані гонки даних, ані стану гонки.%bad

( 3 ) Тут може не бути гонки даних.%good

( 4 ) Тут гарантовано нема ані гонки даних, ані стану гонки.%bad
%enditem


\vspace{5mm}
%beginitem
2. Виберіть усі правильні твердження (якщо такі є).\par
\vspace{2mm}

( 1 ) Потік, запущений конструктором класу jthread, має обов’язково зупинитися, якщо пов’язаний об’єкт надішле команду зупинитися.%bad

( 2 ) Щоб розблокувати заблокований \code{recursive\_mutex}, його завжди достатньо розблокувати один раз.%bad

( 3 ) Якщо два обчислення програми на С++ невдало, але синхронізовані, то на них може виникнути гонка даних.%good

( 4 ) Розблоковувати незаблокований м’ютекс можна.%bad
%enditem


\vspace{5mm}
%beginitem
3. Виберіть усі правильні твердження (якщо такі є).\par
\vspace{2mm}

( 1 ) Вільних від очікувань структур даних не існує.%bad

( 2 ) Операція над структурою даних без блокувань завжди виконається за обмежений час.%bad

( 3 ) Взаємоблокування потоків може виникнути тільки за використання м’ютексів.%bad

( 4 ) Кожна вільна від очікувань структура даних є структурою даних, вільною від блокувань.%good
%enditem


\vspace{5mm}
%beginitem
4. Які проблеми можуть виникнути за concurrent-доступу до даних, що зберігаються у двобічно зв’язаному списку з операціями вставити, додати та ітератором, що забезпечує його обхід? Як їх можна уникнути?

%enditem


\vspace{5mm}
%beginitem
5. Чи можна програмувати мовою Java у стилі декларативного програмування? Відповідь пояснити.
%enditem
}
%end
\end {large}

\vspace{4mm}
%end
\newpage
%begin
\begin {large}

\textbf{01.12.2023 Контрольна робота, варіант  \No { 21 }}\par


{
%beginitem
1. Нехай int var=13; Один потік збільшив змінну var на 3, а інший збільшив її в три рази.

гонка даних (data race), стан гонки (race condition)

Виберіть всі правильні твердження.\par
\vspace{2mm}
( 1 ) Тут гарантовано нема гонки даних, але може бути стан гонки.%bad

( 2 ) Тут гарантовано нема ані гонки даних, ані стану гонки.%bad

( 3 ) Тут може не бути гонки даних.%good

( 4 ) Тут гарантовано є гонка даних.%bad
%enditem


\vspace{5mm}
%beginitem
2. Виберіть усі правильні твердження (якщо такі є).\par
\vspace{2mm}

( 1 ) Щоб розблокувати заблокований \code{shared\_mutex}, його завжди достатньо розблокувати один раз.%good

( 2 ) Об’єкт promise можна скопіювати та роздати копії різним потокам, щоб кожен з них міг передати результат своєї роботи.%bad

( 3 ) Тільки той потік, який обнулив значення семафора, може його збільшити.%bad

( 4 ) М’ютекс має бути записаний на початку критичної секції, яку він забезпечує.%bad
%enditem


\vspace{5mm}
%beginitem
3. Виберіть усі правильні твердження (якщо такі є).\par
\vspace{2mm}

( 1 ) Структура даних без блокувань гарантує завершення виконання довільної операції над цією структурою даних за обмежений час.%bad

( 2 ) Структура даних без блокувань гарантує, що жоден потік не буде очікувати доступу до виконання операцій над цією структурою даних.%bad

( 3 ) Структура даних з блокуваннями може бути такою, що використовує атомарні змінні.%good

( 4 ) Структура даних без блокувань та вільна від блокувань структура даних --- це одне те саме.%bad
%enditem


\vspace{5mm}
%beginitem
4. Які проблеми можуть виникнути за concurrent-доступу до даних, що зберігаються у двобічно зв’язаному списку з операціями вставити, додати та ітератором, що забезпечує його обхід? Як їх можна уникнути?

%enditem


\vspace{5mm}
%beginitem
5. Чи можна програмувати мовою Java у стилі структурного програмування? Відповідь пояснити.
%enditem
}
%end
\end {large}

\vspace{4mm}
%end
\newpage
%begin
\begin {large}

\textbf{01.12.2023 Контрольна робота, варіант  \No { 22 }}\par


{
%beginitem
1. Нехай int var=13; Один потік збільшив змінну var на 3, а інший збільшив її в три рази.

гонка даних (data race), стан гонки (race condition)

Виберіть всі правильні твердження.\par
\vspace{2mm}
( 1 ) Тут може не бути стану гонки.%bad

( 2 ) Тут гарантовано є стан гонки.%good

( 3 ) Тут гарантовано є стан гонки, але гонки даних може не бути.%good

( 4 ) Тут гарантовано є і стан гонки, і гонка даних.%bad
%enditem


\vspace{5mm}
%beginitem
2. Виберіть усі правильні твердження (якщо такі є).\par
\vspace{2mm}

( 1 ) Об’єкт future можна скопіювати та роздати копії різним потокам, щоб вони змогли побачити результат.%bad

( 2 ) Виклик функції \code{this\_thread::sleep\_for(5s)} змушує потік, з якого його здійснили заснути якнайменш на 5 секунд.%good

( 3 ) Довільний потік може розблокувати заблокований м’ютекс.%bad

( 4 ) Текст однопотокової програми мовою С++ однозначно визначає послідовність виконуваних обчислень.%bad
%enditem


\vspace{5mm}
%beginitem
3. Виберіть усі правильні твердження (якщо такі є).\par
\vspace{2mm}

( 1 ) Реалізації різних операцій над структурою даних повинні використовувати різні м’ютекси.%bad

( 2 ) Мова С++20 має стандартну бібліотеку конкурентних структур даних.%bad

( 3 ) За використання структурою даних рівно одного м’ютексу, взаємоблокування на її операціях неможливі.%good

( 4 ) Вільна від блокувань структура даних гарантує, що жоден потік не буде очікувати доступу до виконання операцій над цією структурою даних.%bad
%enditem


\vspace{5mm}
%beginitem
4. Що є точкою входу в новий потік, що його запущено засобами бібліотеки thread? Чи може батьківський потік отримати з нього значення-результат? Чи можна і якщо можна, то як, у новий потік передати посилання на дані батьківського потоку?

%enditem


\vspace{5mm}
%beginitem
5. Що з наступного не підтримує ані мова С++, ані її стандартна бібліотека: каррінг, виклик за необхідністю, замикання, ледачі обчислення, строгі (strict) функції? Відповідь пояснити.
%enditem
}
%end
\end {large}

\vspace{4mm}
%end
\newpage
%begin
\begin {large}

\textbf{01.12.2023 Контрольна робота, варіант  \No { 23 }}\par


{
%beginitem
1. Нехай int var=13; Один потік збільшив змінну var на 3, а інший збільшив її в три рази.

гонка даних (data race), стан гонки (race condition)

Виберіть всі правильні твердження.\par
\vspace{2mm}
( 1 ) Тут гарантовано нема ані гонки даних, ані стану гонки.%bad

( 2 ) Тут гарантовано нема стану гонки, але може бути гонка даних.%bad

( 3 ) Тут може не бути стану гонки.%bad

( 4 ) Тут гарантовано нема стану гонки, але може бути гонка даних.%bad
%enditem


\vspace{5mm}
%beginitem
2. Виберіть усі правильні твердження (якщо такі є).\par
\vspace{2mm}

( 1 ) У програмі може бути не більше однієї критичної секції.%bad

( 2 ) Потік в активному очікуванні споживає ресурс процесора.%good

( 3 ) Досвічений програміст за однопотоковою програмою (текстом мовою С++) завжди може вказати послідовність виконуваних обчислень на конкретних вхідних даних.%bad

( 4 ) Коли програма спить, то вона не виконується.%bad
%enditem


\vspace{5mm}
%beginitem
3. Виберіть усі правильні твердження (якщо такі є).\par
\vspace{2mm}

( 1 ) Мова С++20 має стандартну бібліотеку вільних від блокувань структур даних.%bad

( 2 ) Одночасно використовувати атомарних змінних більше, ніж регістрів процесора, у програмі не можна.%bad

( 3 ) Операція над структурою даних без блокувань завжди виконається за обмежений час.%bad

( 4 ) Операція над вільною від блокувань структурою даних може виконуватись скільки завгодно довго.%good
%enditem


\vspace{5mm}
%beginitem
4. Які бібліотечні засоби С++20 та за якою схемою (кратко) здатні забезпечувати синхронізацію потоків?

%enditem


\vspace{5mm}
%beginitem
5. Чи підходить мова Haskell для програмування в імперативному стилі? Відповідь пояснити.
%enditem
}
%end
\end {large}

\vspace{4mm}
%end
\newpage
%begin
\begin {large}

\textbf{01.12.2023 Контрольна робота, варіант  \No { 24 }}\par


{
%beginitem
1. Нехай int var=13; Один потік збільшив змінну var на 3, а інший збільшив її в три рази.

гонка даних (data race), стан гонки (race condition)

Виберіть всі правильні твердження.\par
\vspace{2mm}
( 1 ) Тут може не бути ані гонки даних, ані стану гонки.%bad

( 2 ) Тут може не бути стану гонки.%bad

( 3 ) Тут може не бути ані гонки даних, ані стану гонки.%bad

( 4 ) Тут може не бути ані гонки даних, ані стану гонки.%bad
%enditem


\vspace{5mm}
%beginitem
2. Виберіть усі правильні твердження (якщо такі є).\par
\vspace{2mm}

( 1 ) Функція \code{sleep\_for} може змусити інший потік заснути.%bad

( 2 ) Чи запустить функція  async виконання в новому потоці можна задавати політикою.%good

( 3 ) Значення, що є результатом виконання функції, запущеної через async, зберігається безпосередньо в об’єкті ф’ючерса (future).%bad

( 4 ) Клас thread має засоби, щоб повернути результат виконання запущеного потоку.%bad
%enditem


\vspace{5mm}
%beginitem
3. Виберіть усі правильні твердження (якщо такі є).\par
\vspace{2mm}

( 1 ) Вільних від очікувань структур даних не існує.%bad

( 2 ) За використання атомарних змінних структура даних є структурою даних без блокувань.%bad

( 3 ) Стандартні контейнери stl безпечно використовувати з кількох потоків.%bad

( 4 ) Реалізація структури даних з блокуваннями може використовувати атомарні змінні.%good
%enditem


\vspace{5mm}
%beginitem
4. Що з точки зору С++20 можна сказати про послідовність обчислень підвиразів у виразі f(a+b)+g(c+d) ? Вважати всі імена коректно означеними.

%enditem


\vspace{5mm}
%beginitem
5. Що з наведеного переліку (якщо таке в переліку є):

\vspace{4mm}
\qquad нормальний порядок редукцій, замикання, каррінг, апплікативний порядок редукцій

\vspace{4mm}
має найслабший логічний зв’язок з більшістю понять (відноситься до іншої парадигми та/або виривається з контексту чи не дуже сумісне)? Відповідь пояснити.

%enditem
}
%end
\end {large}

\vspace{4mm}
%end
\newpage
%begin
\begin {large}

\textbf{01.12.2023 Контрольна робота, варіант  \No { 25 }}\par


{
%beginitem
1. Нехай int var=13; Один потік збільшив змінну var на 3, а інший збільшив її в три рази.

гонка даних (data race), стан гонки (race condition)

Виберіть всі правильні твердження.\par
\vspace{2mm}
( 1 ) Тут гарантовано нема стану гонки, але може бути гонка даних.%bad

( 2 ) Тут може не бути гонки даних.%good

( 3 ) Тут гарантовано нема ані гонки даних, ані стану гонки.%bad

( 4 ) Тут гарантовано є гонка даних.%bad
%enditem


\vspace{5mm}
%beginitem
2. Виберіть усі правильні твердження (якщо такі є).\par
\vspace{2mm}

( 1 ) Значення, що є результатом виконання функції, запущеної через async, зберігається безпосередньо в об’єкті проміза (promise).%bad

( 2 ) Секція є критичною, якщо під час її виконання не можна генерувати винятки.%bad

( 3 ) Блокувати заблокований м’ютекс  \code{mutex} іноді можна.%good

( 4 ) Виклик \code{this\_thread::sleep\_for(5s)} змушує потік проспати рівно 5 секунд.%bad
%enditem


\vspace{5mm}
%beginitem
3. Виберіть усі правильні твердження (якщо такі є).\par
\vspace{2mm}

( 1 ) Структура даних без блокувань гарантує, що жоден потік не буде очікувати доступу до виконання операцій над цією структурою даних.%bad

( 2 ) Операція над вільною від блокувань структурою даних завжди виконається за обмежений час.%bad

( 3 ) Структура даних є структурою даних без блокувань, якщо її реалізація не використовує м’ютекси.%bad

( 4 ) Структура даних може використовувати м’ютекси та атомарні змінні одночасно.%good
%enditem


\vspace{5mm}
%beginitem
4. Чи завжди наслідком стану гонки (race condition) є невизначена поведінка програми? Відповідь пояснити.

%enditem


\vspace{5mm}
%beginitem
5. Що з наведеного переліку (якщо таке в переліку є):

\vspace{4mm}
\qquad функції вищого порядку, анонімні функції, функції першого порядку, замикання

\vspace{4mm}
має найслабший логічний зв’язок з більшістю понять (відноситься до іншої парадигми та/або виривається з контексту чи не дуже сумісне)? Відповідь пояснити.

%enditem
}
%end
\end {large}

\vspace{4mm}
%end
\newpage
%begin
\begin {large}

\textbf{01.12.2023 Контрольна робота, варіант  \No { 26 }}\par


{
%beginitem
1. Нехай int var=13; Один потік збільшив змінну var на 3, а інший збільшив її в три рази.

гонка даних (data race), стан гонки (race condition)

Виберіть всі правильні твердження.\par
\vspace{2mm}
( 1 ) Тут гарантовано є і стан гонки, і гонка даних.%bad

( 2 ) Тут гарантовано є стан гонки, але гонки даних може не бути.%good

( 3 ) Тут може не бути гонки даних.%good

( 4 ) Тут гарантовано є стан гонки.%good
%enditem


\vspace{5mm}
%beginitem
2. Виберіть усі правильні твердження (якщо такі є).\par
\vspace{2mm}

( 1 ) Функція async запускає виконання в новому потоці.%bad

( 2 ) Тільки від компілятора залежить, чи запустить функція async виконання в новому потоці.%bad

( 3 ) Розблоковувати незаблокований м’ютекс не можна.%good

( 4 ) Критична секція це неперервна ділянка коду програми.%bad
%enditem


\vspace{5mm}
%beginitem
3. Виберіть усі правильні твердження (якщо такі є).\par
\vspace{2mm}

( 1 ) Структура даних без блокувань гарантує завершення виконання довільної операції над цією структурою даних за обмежений час.%bad

( 2 ) Якщо весь код структури даних утворює одну критичну секцію, то така структура даних може використовуватись в конкурентній програмі, але по-справжньому конкурентною не є.%good

( 3 ) Мова С++20 має стандартну бібліотеку вільних від блокувань структур даних.%bad

( 4 ) Одночасно використовувати атомарних змінних більше, ніж регістрів процесора, у програмі не можна.%bad
%enditem


\vspace{5mm}
%beginitem
4. Які бібліотечні засоби С++20 можна використати, що забезпечити передачу результатів обчислення від одного потоку до іншого? Яка схема їх використання?

%enditem


\vspace{5mm}
%beginitem
5. Що з наведеного переліку (якщо таке в переліку є):

\vspace{4mm}
\qquad блок, стан пам’яті, побічний ефект, замикання

\vspace{4mm}
має найслабший логічний зв’язок з більшістю понять (відноситься до іншої парадигми та/або виривається з контексту чи не дуже сумісне)? Відповідь пояснити.

%enditem
}
%end
\end {large}

\vspace{4mm}
%end
\newpage
\end{document}

